% English translation of checked_pages/esq_119_037.tex; source: cote 119, site 37, PDF page 38, typescript page 35.
% AI-translated by GPT-6 Astra (OpenAI, ChatGPT Pro mode).
\scanpage{119/37}
\mspage{35}

situations, for example for the compact moduli space
$M^{\wedge}_{g,\nu}$ with its canonical ``stratification at infinity'', given
by the Mumford-Deligne normal crossings divisor. Another
difficulty, probably less serious, is that the so-called
moduli ``space'' is in fact a \uline{multiplicity} --- technically,
this is expressed above all by the need to replace
the index set $I$ for the strata by an (essentially finite)
index \uline{category}, in this case that of
``MD graphs'', which ``parametrize'' the possible ``combinatorial structures''
of a stable curve of type $(g,\nu)$. That said, I can
affirm that the general theory of dévissage, developed especially
under the pressure of the needs of \uline{this} cause, has indeed proved
a valuable guide, leading to a progressive understanding,
with unfailing coherence, of certain essential aspects
of the Teichmüller tower (that is, essentially
of the ``structure at infinity'' of the ordinary Teichmüller groups).
It is this approach that finally led me, in the following
months, toward the principle of a purely combinatorial
construction of the tower of Teichmüller groupoids, in the spirit
sketched \unsure{.}above\tnote{TN-037-01}{The scan has a doubtful dot between ``plus'' and ``haut'' in ``plus haut''; it is retained beside the corresponding English word.} (cf. par. 2).

Another satisfactory test of coherence comes from the
``topossic'' viewpoint.\tnote{TN-037-02}{``Topossic'' renders the author's adjective ``topossique'', referring here to the context of topoi.} Indeed, since my interest in moduli
multiplicities stems above all from their algebraic-geometric
and arithmetic meaning, it is the \uline{algebraic} moduli
multiplicities, over the absolute base field $\mathbb{Q}$, that have
primarily interested me, and a dévissage ``at infinity'' of their
geometric fundamental groups (i.e. the \uline{profinite} Teichmüller
groups) compatible with the natural actions of
$\mathrm{I}\!\Gamma = \mathrm{Gal}(\bar{\mathbb{Q}}/\mathbb{Q})$. This seemed from the outset to exclude the possibility of
referring to a hypothetical theory of dévissage of stratified
structures in a context of ``tame topology'' (or even of
ordinary topology, after a fashion), except as an
entirely heuristic guiding thread. The question then arose of translating,
into the context of the topoi (in this case the étale topoi)

\margin{37}
